% =============================================================================
%  text/RESULTS/main-results.tex
%  Main results: demonstrates a regression-table input and figure reference.
%  主要结果示例：引用回归表与图片。
% =============================================================================

Table~\ref{tab:main-promotion} reports estimates of
equation~\eqref{eq:promotion} for the period in which the ruler's authority
was consolidated. Consistent with Hypothesis~\ref{hyp:promotion}, the
coefficient on network centrality is positive and statistically significant
across specifications, including the full set of controls.

\begin{figure}[htbp]
  \centering
  \includegraphics[width=0.8\textwidth]{example-figure.pdf}
  \caption{Marginal Effects of Network Centrality on Promotion (Placeholder)}
  \label{fig:marginal-effects}
\end{figure}

% 中文对照（可选）
% 表 B1 报告权力巩固时期晋升方程的估计结果：网络中心性的系数为正且显著，
% 与假设 1 一致。图 2 展示了核心变量对晋升的边际效应（示意图，请替换）。
Figure~\ref{fig:marginal-effects} visualizes the marginal effects.

Table~\ref{tab:main-purge} reports estimates of equation~\eqref{eq:purge}
for the period in which the ruler's authority was challenged. Consistent with
Hypothesis~\ref{hyp:suppression}, the coefficient on network centrality is
positive and statistically significant: precisely the elites who were
rewarded during the consolidation period became targets of suppression when
the ruler's authority wavered.

% 表 B2 报告权力受挑战时期清洗方程的估计结果：网络中心性系数为正且显著，
% 与假设 2 一致——权力巩固时期受奖赏的精英在权威动摇时反而成为清洗对象。
These two sets of results jointly support the dynamic framework developed in
Section~\ref{sec:elite-ruler}.
